\chapter{总结与展望}
\label{chap:conclusion}

本章对全文工作进行总结，概括本文在 Agent 自主任务调度方法、LLM 分层信息整合策略和工程系统实现方面取得的主要成果，并在此基础上展望系统后续可以继续拓展的方向。

\section{全文总结}
\label{sec:conclusion-full-text-summary}

本文围绕建筑幕墙韧性评估中的现实痛点，设计并实现了一套基于 AI Agent 的建筑幕墙韧性评估软件系统。有别于围绕单个检测模型展开的研究路线，本文重点解决现实工程中更关键的链路问题：如何把项目级图像资产、多类原子检测能力、人工复核信息和项目级报告生成组织为一个完整、可追踪、可部署的软件原型系统，并通过 Agent 提升任务调度和信息整合能力。

全文首先分析了建筑幕墙巡检从人工经验向智能评估转型过程中面临的痛点，指出现有单点检测方法难以解决项目级图像管理、检测任务组织和建筑级报告生成问题。随后，本文梳理了建筑立面检测、计算机视觉和 AI Agent 相关研究，说明 Agent 工具调用和多阶段工作流为幕墙评估提供了新的技术机会。在问题建模基础上，本文提出 Agent 自主任务调度方法和 LLM 分层信息整合策略，将分类 Agent、图像级总结 Agent 和项目级报告 Agent 分别嵌入检测前、检测后和项目报告阶段。最后，本文完成了包含 Web 前端、Core API、Core BIZ、Classification、Reasoning、Summary、ICW Common、MySQL、Redis、MinIO、RocketMQ 和 Docker Compose 部署配置的系统原型，并通过系统验证展示了完整业务闭环。

\section{主要成果}
\label{sec:conclusion-main-results}

本文取得的主要成果包括以下三个方面。

第一，提出并实现了面向多材质幕墙图像的 Agent 自主任务调度方法，使系统能够根据图像内容动态决定需要执行的检测能力。该方法将分类 Agent 的语义判断约束为有限任务代码集合，并由 Core BIZ 统一转换为检测主任务和检测子任务，从而把开放的图像理解能力转化为可治理、可追踪的工程任务调度过程。该成果解决了多材质幕墙图像中“哪些图像需要执行哪些检测能力”的问题，使检测流程更符合幕墙材料差异和实际业务需求。

第二，提出并实现了 LLM 分层信息整合策略。系统先保存原子检测结构化结果，再生成图像级总结，最后生成项目级报告。该策略将“单张图像说明什么”和“整个项目说明什么”拆分为不同语义层次，将大语言模型从简单文本生成工具扩展为“路由、压缩、汇总”的智能节点，减少了多图像、多任务、多区域信息直接汇总输入 LLM 模型带来的上下文过载和重点淹没问题，使报告生成过程更稳定、更可解释，也更符合土木工程领域用户从局部缺陷到整体评估的认知路径。

第三，构建了完整工程闭环。系统通过微服务架构、gRPC 协议、异步 Worker、回调接口、数据库状态机、对象存储、消息队列和 WebSocket 推送，使 Agent 能力从一次性调用变成可运行、可追踪、可恢复的业务流程。前端页面能够展示项目状态、图像资产、检测进度、检测产物、人工复核和项目报告，后端服务能够完成任务创建、状态推进、产物校验、失败重试和部署运行。本成果体现了本系统作为软件工程专业毕业设计的工作完整性。

\section{未来工作展望}
\label{sec:conclusion-future-work-outlook}

基于本文已经完成的软件原型系统，未来工作可以从以下三个方面继续推进。

第一，可以进一步扩充建筑幕墙真实场景数据集，对金属锈蚀、石材裂缝、石材污渍、玻璃平整度和玻璃爆裂五类原子检测能力进行更系统的精度评估和模型优化。本文重点在于系统链路和 Agent 工程化方法，原子检测能力主要作为工具化能力接入系统。未来若能够结合更多真实建筑、不同拍摄角度、不同光照环境和不同幕墙材料的数据进行模型训练与评估，系统的检测结果将具有更强的工程适应性。

第二，可以继续增强系统工程能力。本文已经完成了容器化部署、日志采集、异步任务、消息事件、WebSocket 推送和失败重试等基础能力，后续可进一步加入分布式 Worker、任务优先级调度、模型版本管理、监控告警机制和更细粒度的权限控制。这些能力有助于系统在软件原型系统的基础上进一步走向更稳定的工程化应用。

第三，可以探索系统与 BIM 平台、无人机巡检平台和建筑运维系统的对接。建筑幕墙韧性评估不仅需要识别图像中的缺陷，还需要进一步理解缺陷在建筑空间中的位置、发展趋势和维护优先级。未来若能够将图像检测结果映射到 BIM 构件、立面区域和历史维修记录中，系统将能够从图像级评估进一步扩展到建筑全生命周期的数字运维与智能管理。

综上所述，本文验证了 AI Agent 在建筑幕墙韧性评估中的应用价值。Agent 不是替代工程系统，而是通过自主任务调度和分层信息整合能力补足传统工程系统在语义判断和报告生成方面的不足。本文工作为建筑幕墙智能运维提供了一条可落地的技术路线，也为 AI Agent 在垂直工程软件中的应用提供了参考。
